% Unit 038 terjemahan Bahasa Indonesia dari chapter3.tex baris 1061--1291.
% Sumber: Methods of Algebra, Volume 2, commit
% 9a5803ff77dd3257484cb177f851a73770a59dd3 (CC BY 4.0).
% stable-unit-id: o014.aljabr2.chapter3.mapping-cone-and-long-exact-sequences
% Cakupan: Bagian 3.7 lengkap.

% segment-id: o014.li2.chapter3.mapping-cone-and-long-exact-sequences.h001
\section{Kerucut pemetaan: barisan eksak panjang}\label{sec:cone-vs-long-exact-sequence}
\hypertarget{o014.aljabr2.chapter3.mapping-cone-and-long-exact-sequences}{ }

% segment-id: o014.li2.chapter3.mapping-cone-and-long-exact-sequences.p001
Barisan eksak panjang merupakan alat utama dalam aljabar homologis. Bagian ini
bertujuan membandingkan tiga cara membangun barisan eksak panjang dari barisan
eksak pendek. Terlepas dari beberapa perbedaan tanda, ketiganya memberikan
hasil yang sama. Dalam bagian ini, $\mathcal{A}$ diasumsikan sebagai kategori
abelian.

% segment-id: o014.li2.chapter3.mapping-cone-and-long-exact-sequences.pr001
\begin{proposisi}\label{prop:cone-triangle-ses}
	Misalkan $f:X\to Y$ merupakan morfisme dalam $\cate{C}(\mathcal{A})$.
	Ambil $\alpha(f)$ dan $\beta(f)$ sebagaimana dalam
	Definisi~\ref{def:cone-triangle}. Maka
% segment-id: o014.li2.chapter3.mapping-cone-and-long-exact-sequences.q001
	\[ 0 \to Y \xrightarrow{\alpha(f)} \Cone(f) \xrightarrow{\beta(f)} X[1] \to 0 \]
	merupakan barisan eksak dalam $\cate{C}(\mathcal{A})$.
\end{proposisi}

% segment-id: o014.li2.chapter3.mapping-cone-and-long-exact-sequences.pf001
\begin{bukti}
	Ingat kembali definisi $\alpha(f)$ dan $\beta(f)$. Karena keeksakan dapat
	diperiksa derajat demi derajat (Proposisi~\ref{prop:abelian-cat-cplx}),
	pernyataan tersebut cukup dibuktikan melalui keeksakan
% segment-id: o014.li2.chapter3.mapping-cone-and-long-exact-sequences.q002
	\[ 0 \to Y^n \xrightarrow{(0, \identity)} X^{n+1} \oplus Y^n \xrightarrow{\text{proyeksi}} X^{n+1} \to 0, \]
	untuk setiap $n\in\Z$. Inilah isi Proposisi~\ref{prop:biproduct-ses}.
\end{bukti}

% segment-id: o014.li2.chapter3.mapping-cone-and-long-exact-sequences.p002
Diberikan morfisme $f:X\to Y$ dalam $\cate{C}(\mathcal{A})$, barisan eksak
pendek pada Proposisi~\ref{prop:cone-triangle-ses}, bersama dengan
Proposisi~\ref{prop:long-exact-sequence-ses}, memberikan barisan eksak panjang
berikut dalam $\mathcal{A}$:
% segment-id: o014.li2.chapter3.mapping-cone-and-long-exact-sequences.q003
\begin{equation}\label{eqn:cone-long-exact-sequence}
	\cdots \to \Hm^n(Y) \xrightarrow{\Hm^n(\alpha(f))} \Hm^n(\Cone(f)) \xrightarrow{\Hm^n(\beta(f))} \underbracket{\Hm^n(X[1])}_{= \Hm^{n+1}(X)} \xrightarrow{\xi^n} \Hm^{n+1}(Y) \to \cdots
\end{equation}
Morfisme yang ditandai dengan $\xi^n$ ialah morfisme penghubung yang
diinduksi oleh barisan eksak pendek tersebut. Korolari~\ref{prop:cone-connecting}
akan membuktikan bahwa $\xi^n$ tidak lain ialah
$\Hm^{n+1}(f):\Hm^{n+1}(X)\to\Hm^{n+1}(Y)$. Dengan demikian, setiap bagian
dari barisan eksak panjang kerucut pemetaan
\eqref{eqn:cone-long-exact-sequence} bersifat eksplisit dan berasal dari
$f$, $\alpha(f)$, dan $\beta(f)$.

% segment-id: o014.li2.chapter3.mapping-cone-and-long-exact-sequences.p003
Hasil berikut menunjukkan arah sebaliknya: dari sembarang barisan eksak pendek
dalam $\cate{C}(\mathcal{A})$, terdapat pula hubungan alami dengan kerucut
pemetaan. Hubungan ini dapat dibentuk dengan dua cara.

% segment-id: o014.li2.chapter3.mapping-cone-and-long-exact-sequences.le001
\begin{lema}\label{prop:cone-qis}
	Diberikan barisan eksak pendek
	$0\to X\xrightarrow{f}Y\xrightarrow{g}Z\to0$ dalam
	$\cate{C}(\mathcal{A})$, kedua morfisme
% segment-id: o014.li2.chapter3.mapping-cone-and-long-exact-sequences.q004
	\[ \Phi := (0, g): \Cone(f) \to Z, \quad \Phi' := (f[1], 0): X[1] \to \Cone(g) \]
	merupakan kuasi-isomorfisme antarkompleks dan mempunyai sifat-sifat berikut:
% segment-id: o014.li2.chapter3.mapping-cone-and-long-exact-sequences.l001
	\begin{enumerate}[(i)]
% segment-id: o014.li2.chapter3.mapping-cone-and-long-exact-sequences.i001
		\item $\Phi \circ \alpha(f) = g$ dan $\beta(g) \circ \Phi' = f[1]$;
% segment-id: o014.li2.chapter3.mapping-cone-and-long-exact-sequences.i002
		\item $\alpha(g)\Phi + \Phi'\beta(f):\Cone(f)\to\Cone(g)$ secara
		kanonik homotopik nol.
	\end{enumerate}
\end{lema}

% segment-id: o014.li2.chapter3.mapping-cone-and-long-exact-sequences.pf002
\begin{bukti}
	Perhatikan diagram komutatif berikut dalam $\cate{C}(\mathcal{A})$:
% segment-id: o014.li2.chapter3.mapping-cone-and-long-exact-sequences.g001
	\[\begin{tikzcd}[column sep=small]
		0 \arrow[r] & X \arrow[r, "\identity_X"] \arrow[d, "{\identity_X}"'] & X \arrow[r] \arrow[d, "f"] & 0 \arrow[r] \arrow[d] & 0 \\
		0 \arrow[r] & X \arrow[r, "f"'] & Y \arrow[r, "g"'] & Z \arrow[r] & 0
	\end{tikzcd} \quad \begin{tikzcd}[column sep=small]
		0 \arrow[r] & X \arrow[d] \arrow[r, "f"] & Y \arrow[d, "g"'] \arrow[r, "g"] & Z \arrow[d, "\identity_Z"] \arrow[r] & 0 \\
		0 \arrow[r] & 0 \arrow[r] & Z \arrow[r, "\identity_Z"'] & Z \arrow[r] & 0
	\end{tikzcd}\]
	Setiap baris eksak. Oleh karena itu, funktorialitas kerucut pemetaan
	(Proposisi~\ref{prop:cone-functorial}) memberikan morfisme-morfisme dalam
	$\cate{C}(\mathcal{A})$ yang tersusun sebagai
% segment-id: o014.li2.chapter3.mapping-cone-and-long-exact-sequences.g002
	\begin{equation}\label{eqn:cone-qis-aux}\begin{tikzcd}[row sep=tiny]
		0 \arrow[r] & \Cone\left( \identity_X \right) \arrow[r] & \Cone\left( X \xrightarrow{f} Y \right) \arrow[r] & \Cone\left( 0 \to Z \right) \arrow[r] & 0 \\
		0 \arrow[r] & \Cone(X \to 0) \arrow[r] & \Cone\left( Y \xrightarrow{g} Z \right) \arrow[r] & \Cone\left( \identity_Z \right) \arrow[r] & 0
	\end{tikzcd}\end{equation}
	Dengan memeriksa derajat demi derajat serta menggunakan diagram-diagram
	komutatif di atas dan definisi kerucut pemetaan, tampak bahwa setiap baris
	dalam~\eqref{eqn:cone-qis-aux} eksak. Selain itu, mudah dilihat bahwa
	$\Cone(f)\to\Cone(0\to Z)\simeq Z$ (atau
	$X[1]\simeq\Cone(X\to0)\to\Cone(g)$) tepat merupakan $\Phi$ (atau
	$\Phi'$). Setelah diketahui bahwa keduanya merupakan morfisme kompleks,
	persamaan-persamaan dalam (i) langsung terlihat.

	Dengan menerapkan Proposisi~\ref{prop:long-exact-sequence-ses}
	pada~\eqref{eqn:cone-qis-aux}, untuk setiap $n\in\Z$ diperoleh barisan eksak
% segment-id: o014.li2.chapter3.mapping-cone-and-long-exact-sequences.q005
	\begin{gather*}
		\Hm^n\left( \Cone(\identity_X) \right) \to \Hm^n \left( \Cone(f) \right) \xrightarrow{\Hm^n(\Phi)} \Hm^n(Z) \to \Hm^{n+1} \left( \Cone(\identity_X) \right), \\
		\Hm^{n-1}\left( \Cone(\identity_Z) \right) \to \Hm^n(X[1]) \xrightarrow{\Hm^n(\Phi')} \Hm^n\left( \Cone(g) \right) \to \Hm^n\left( \Cone(\identity_Z) \right).
	\end{gather*}
	Lema~\ref{prop:cone-homotopy} (i) dan
	Proposisi~\ref{prop:null-homotopic-cohomology} menunjukkan bahwa suku-suku
	ujung kiri dan kanan semuanya $0$. Jadi, $\Phi$ dan $\Phi'$ merupakan
	kuasi-isomorfisme.

	Terakhir, kita verifikasi (ii). Dengan menggunakan
	$\Cone(f)^n=X^{n+1}\oplus Y^n$ dan
	$\Cone(g)^n=Y^{n+1}\oplus Z^n$, tuliskan morfisme-morfisme tersebut sebagai
	matriks:
% segment-id: o014.li2.chapter3.mapping-cone-and-long-exact-sequences.g003
	\begin{gather*}\begin{tikzcd}[row sep=tiny, column sep=large, ampersand replacement=\&]
		\& Z \arrow[rd, "\alpha(g)" description, "{\bigl(\begin{smallmatrix} 0 \\ \identity \end{smallmatrix} \bigr)}" inner sep=0.6em] \& \\
		\Cone(f) \arrow[ru, "\Phi" description, "{\bigl( 0 \; g \bigr)}" inner sep=0.6em] \arrow[rd, "\beta(f)" description, "{\big( \identity \; 0 \bigr)}"' inner sep=0.6em] \& \& \Cone(g) \\
		\& X[1] \arrow[ru, "{\Phi'}" description, "{\bigl( \begin{smallmatrix} f[1] \\ 0 \end{smallmatrix} \bigr)}"' inner sep=0.6em] \&
	\end{tikzcd}, \quad
		\alpha(g) \Phi + \Phi' \beta(f) = \begin{pmatrix} f[1] & 0 \\ 0 & g  \end{pmatrix}.
	\end{gather*}
	Di sisi lain, ambil $s=(s^n)_n\in
	\Hom^{-1}\left(\Cone(f),\Cone(g)\right)$ dengan
% segment-id: o014.li2.chapter3.mapping-cone-and-long-exact-sequences.q006
	\[ s^n := \begin{pmatrix}
		0 & \identity_{Y^n} \\
		0 & 0
	\end{pmatrix}: X^{n+1} \oplus Y^n \to Y^n \oplus Z^{n-1} , \]
	yang juga merupakan satu-satunya pilihan yang wajar. Perhitungan matriks
	langsung menunjukkan bahwa
	$d_{\Cone(g)}^{n-1}s^n+s^{n+1}d_{\Cone(f)}^n$ sama dengan
% segment-id: o014.li2.chapter3.mapping-cone-and-long-exact-sequences.q007
	\[\begin{pmatrix}
		-d_Y^n & 0 \\
		g^n & d_Z^{n-1}
	\end{pmatrix} \begin{pmatrix}
		0 & \identity_{Y^n} \\
		0 & 0
	\end{pmatrix} + \begin{pmatrix}
		0 & \identity_{Y^{n+1}} \\
		0 & 0
	\end{pmatrix} \begin{pmatrix}
		-d_X^{n+1} & 0 \\
		f^{n+1} & d_Y^n
	\end{pmatrix} = \begin{pmatrix}
		f^{n+1} & 0 \\
		0 & g^n
	\end{pmatrix}.\]
	Jadi, $\alpha(g)\Phi+\Phi'\beta(f)$ memang homotopik nol.
\end{bukti}

% segment-id: o014.li2.chapter3.mapping-cone-and-long-exact-sequences.p004
Sebagai rangkuman, dari barisan eksak pendek
$0\to X\xrightarrow{f}Y\xrightarrow{g}Z\to0$ dalam
$\cate{C}(\mathcal{A})$, terdapat tiga cara untuk memperoleh barisan eksak
panjang berikut dalam $\mathcal{A}$:
% segment-id: o014.li2.chapter3.mapping-cone-and-long-exact-sequences.g004
\[\begin{tikzcd}[row sep=large]
	\cdots \arrow[r] & \Hm^{n-1}(Y) \arrow[r] \arrow[phantom, d, ""{coordinate, name=A}] & \Hm^{n-1}(Z) \arrow[dll, rounded corners, to path={
		-- ([xshift=2ex]\tikztostart.east)
		|- (A) [near end]\tikztonodes
		-| ([xshift=-2ex]\tikztotarget.west)
		-- (\tikztotarget)}] \\
	\Hm^n(X) \arrow[r] & \Hm^n(Y) \arrow[r] \arrow[phantom, d, ""{coordinate, name=B}] & \Hm^n(Z) \arrow[dll, rounded corners, to path={
		-- ([xshift=2ex]\tikztostart.east)
		|- (B) [near end]\tikztonodes
		-| ([xshift=-2ex]\tikztotarget.west)
		-- (\tikztotarget)}] \\
	\Hm^{n+1}(X) \arrow[r] & \Hm^{n+1}(Y) \arrow[r] & \cdots .
\end{tikzcd}\]
% segment-id: o014.li2.chapter3.mapping-cone-and-long-exact-sequences.l002
\begin{enumerate}[(A)]
% segment-id: o014.li2.chapter3.mapping-cone-and-long-exact-sequences.i003
	\item Dengan menerapkan Proposisi~\ref{prop:long-exact-sequence-ses}
	secara langsung, diperoleh barisan eksak panjang
% segment-id: o014.li2.chapter3.mapping-cone-and-long-exact-sequences.g005
	\[\begin{tikzcd}
		\cdots \arrow[r] & \Hm^n(Y) \arrow[r, "{\Hm^n(g)}" inner sep=0.6em] & \Hm^n(Z) \arrow[r, "{\delta^n}" inner sep=0.6em] & \Hm^{n+1}(X) \arrow[r, "{\Hm^{n+1}(f)}" inner sep=0.6em] & \Hm^{n+1}(Y) \arrow[r] & \cdots
	\end{tikzcd}\]
% segment-id: o014.li2.chapter3.mapping-cone-and-long-exact-sequences.i004
	\item Gunakan kuasi-isomorfisme $\Phi:\Cone(f)\to Z$ dari
	Lema~\ref{prop:cone-qis}, bersama dengan
	\eqref{eqn:cone-long-exact-sequence}, untuk memperoleh diagram komutatif
% segment-id: o014.li2.chapter3.mapping-cone-and-long-exact-sequences.g006
	\[\begin{tikzcd}[column sep=scriptsize]
		\cdots \arrow[r] & \Hm^n(Y) \arrow[r, "{\Hm^n(\alpha(f))}" inner sep=0.6em] \arrow[equal, d] & \Hm^n\left( \Cone(f) \right) \arrow[r, "{\Hm^n(\beta(f))}" inner sep=0.6em] \arrow[d, "{\Hm^n(\Phi)}"] & \Hm^n(X[1]) \arrow[r] \arrow[equal, d] & \Hm^{n+1}(Y) \arrow[r] \arrow[equal, d] & \cdots \\
		\cdots \arrow[r] & \Hm^n(Y) \arrow[r, "{\Hm^n(g)}"' inner sep=0.6em] & \Hm^n(Z) \arrow[r, "{\eta^n}"' inner sep=0.6em] & \Hm^{n+1}(X) \arrow[r, "{\xi^n}"' inner sep=0.6em] & \Hm^{n+1}(Y) \arrow[r] & \cdots
	\end{tikzcd}\]
	dengan
% segment-id: o014.li2.chapter3.mapping-cone-and-long-exact-sequences.l003
	\begin{compactitem}
% segment-id: o014.li2.chapter3.mapping-cone-and-long-exact-sequences.i005
		\item $\eta^n := \Hm^n(\beta(f)) \Hm^n(\Phi)^{-1}$;
% segment-id: o014.li2.chapter3.mapping-cone-and-long-exact-sequences.i006
		\item $\xi^n$ ialah morfisme penghubung dalam
		\eqref{eqn:cone-long-exact-sequence}, yang berbentuk
		\[
			\Hm^{n+1}(X)=\Hm^n(X[1])\to\Hm^{n+1}(Y).
		\]
	\end{compactitem}
	Karena baris pertama diketahui eksak, baris kedua juga eksak.

% segment-id: o014.li2.chapter3.mapping-cone-and-long-exact-sequences.i007
	\item Lakukan hal yang sama, tetapi kini gunakan kuasi-isomorfisme
	$\Phi':X[1]\to\Cone(g)$ dari Lema~\ref{prop:cone-qis}. Diperoleh diagram
	komutatif
% segment-id: o014.li2.chapter3.mapping-cone-and-long-exact-sequences.g007
	\[\begin{tikzcd}[column sep=scriptsize]
		\cdots \arrow[r] & \Hm^n(Z) \arrow[r, "{\Hm^n(\alpha(g))}" inner sep=0.6em] & \Hm^n\left( \Cone(g) \right) \arrow[r, "{\Hm^n(\beta(g))}" inner sep=0.6em] & \Hm^n(Y[1]) \arrow[r] & \Hm^{n+1}(Z) \arrow[r] & \cdots \\
		\cdots \arrow[r] & \Hm^n(Z) \arrow[r, "{(\eta')^n}"' inner sep=0.6em] \arrow[equal, u] & \Hm^n(X[1]) \arrow[r, "{\Hm^n(f[1])}"' inner sep=0.6em] \arrow[u, "{\Hm^n(\Phi')}"', "\simeq"] & \Hm^{n+1}(Y) \arrow[r, "{(\xi')^n}"' inner sep=0.6em] \arrow[equal, u] & \Hm^{n+1}(Z) \arrow[r] \arrow[equal, u] & \cdots
	\end{tikzcd}\]
	dengan
% segment-id: o014.li2.chapter3.mapping-cone-and-long-exact-sequences.l004
	\begin{compactitem}
% segment-id: o014.li2.chapter3.mapping-cone-and-long-exact-sequences.i008
		\item $(\eta')^n := \Hm^n(\Phi')^{-1} \Hm^n(\alpha(g))$;
% segment-id: o014.li2.chapter3.mapping-cone-and-long-exact-sequences.i009
		\item $(\xi')^n$ berasal dari morfisme penghubung
		$\Hm^{n+1}(Y)=\Hm^n(Y[1])\to\Hm^{n+1}(Z)$
		dalam~\eqref{eqn:cone-long-exact-sequence}, dengan $g$ menggantikan $f$.
	\end{compactitem}
	Karena baris pertama diketahui eksak, baris kedua juga eksak.
\end{enumerate}

% segment-id: o014.li2.chapter3.mapping-cone-and-long-exact-sequences.p005
Apa persamaan dan perbedaan di antara ketiga barisan eksak panjang ini? Kita
nyatakan terlebih dahulu kesimpulannya.

% segment-id: o014.li2.chapter3.mapping-cone-and-long-exact-sequences.l005
\begin{itemize}
% segment-id: o014.li2.chapter3.mapping-cone-and-long-exact-sequences.i010
	\item Konstruksi (A) dan (B) hanya berbeda dalam beberapa tanda negatif:
	Proposisi~\ref{prop:connecting-eta-delta} akan menunjukkan bahwa
	$\eta^n=-\delta^n$, sedangkan Korolari~\ref{prop:cone-connecting}
	menyiratkan $\xi^n=\Hm^{n+1}(f)$.
% segment-id: o014.li2.chapter3.mapping-cone-and-long-exact-sequences.i011
	\item Konstruksi (A) dan (C) memberikan hasil yang sama:
	Korolari~\ref{prop:connecting-eta-delta-2}, berdasarkan hasil untuk (B),
	akan menunjukkan bahwa $(\eta')^n=\delta^n$, sedangkan
	Korolari~\ref{prop:cone-connecting} menyiratkan
	$(\xi')^n=\Hm^{n+1}(g)$.
% segment-id: o014.li2.chapter3.mapping-cone-and-long-exact-sequences.i012
	\item Dengan demikian, (B) dan (C) hanya berbeda tanda pada morfisme
	penghubung $\Hm^n(Z)\to\Hm^{n+1}(X)$. Hal ini dapat dijelaskan melalui
	aksioma rotasi (TR3) bagi kategori bertriangulasi; lihat
	Proposisi~\sourcecrossref{prop:ses-vs-triangle}{pada pembahasan kategori
	turunan di Bab~4}.
\end{itemize}

% segment-id: o014.li2.chapter3.mapping-cone-and-long-exact-sequences.p006
Selanjutnya, kita buktikan hubungan-hubungan tersebut.

% segment-id: o014.li2.chapter3.mapping-cone-and-long-exact-sequences.pr002
\begin{proposisi}\label{prop:connecting-eta-delta}
	Dalam situasi di atas, $\eta^n=-\delta^n$ untuk setiap $n\in\Z$.
\end{proposisi}

% segment-id: o014.li2.chapter3.mapping-cone-and-long-exact-sequences.pf003
\begin{bukti}
	Argumen berikut diambil dari \cite[Proposisi~12.3.6]{KS06}. Karena
	argumen ini cukup rumit, pembaca dianjurkan terlebih dahulu mencoba kasus
	konkret $\mathcal{A}=R\dcate{Mod}$, dengan $R$ gelanggang; metode pengejaran
	diagram dalam \cite[\S 6.8]{Li1} memberikan bukti langsung untuk kasus
	tersebut.

	Untuk kategori abelian umum $\mathcal{A}$, pertama-tama misalkan
	$n\in\Z$ dan ambil produk serat
% segment-id: o014.li2.chapter3.mapping-cone-and-long-exact-sequences.q008
	\[ S := \Coker(d_Y^{n-1}) \dtimes{\Coker(d_Z^{n-1})} \Hm^n(Z). \]
	Bagaimana morfisme penghubung $\delta^n$ dikonstruksi? Morfisme ini
	diperoleh dengan menerapkan~\eqref{eqn:snake-conn} pada diagram dalam
	\eqref{eqn:long-exact-sequence-ses-aux},
% segment-id: o014.li2.chapter3.mapping-cone-and-long-exact-sequences.g008
	\[\begin{tikzcd}
		& \Coker(d_X^{n-1}) \arrow[r] \arrow[d] & \Coker(d_Y^{n-1}) \arrow[r] \arrow[d] & \Coker(d_Z^{n-1}) \arrow[d] \arrow[r] & 0 \\
		0 \arrow[r] & \Ker(d_X^{n+1}) \arrow[r] & \Ker(d_Y^{n+1}) \arrow[r] & \Ker(d_Z^{n+1}) &
	\end{tikzcd}\]
	yang di sini disesuaikan menjadi bentuk
% segment-id: o014.li2.chapter3.mapping-cone-and-long-exact-sequences.g009
	\begin{equation}\label{eqn:connecting-eta-delta-0}\begin{tikzcd}
		& S \arrow[twoheadrightarrow, r, "w"] \arrow[d, "v"] \arrow[ldd, bend right, "u"'] \arrow[rd, phantom, "\Box" description] & \Hm^n(Z) \arrow[hookrightarrow, d] \\
		& \Coker(d_Y^{n-1}) \arrow[twoheadrightarrow, r] \arrow[d, "a"] & \Coker(d_Z^{n-1}) \arrow[d] \\
		\Ker(d_X^{n+1}) \arrow[hookrightarrow, r, "b"'] \arrow[twoheadrightarrow, d] & \Ker(d_Y^{n+1}) \arrow[r, "c"'] & \Ker(d_Z^{n+1}) \\
		\Hm^{n+1}(X) & &
	\end{tikzcd}\end{equation}
	Di sini semua baris dan kolom eksak, panah $a$ diinduksi oleh $d_Y^n$,
	panah $b$ diinduksi oleh $f^{n+1}$, sedangkan keberadaan dan ketunggalan
	panah $u$ mengikuti dari fakta bahwa komutativitas bagian kanan menyiratkan
	$cav=0$. Dengan meninjau secara saksama konstruksi
	dalam~\eqref{eqn:snake-conn}, tampak bahwa diagram
% segment-id: o014.li2.chapter3.mapping-cone-and-long-exact-sequences.g010
	\begin{equation}\label{eqn:connecting-eta-delta-0.5}\begin{tikzcd}[row sep=tiny]
		& \Hm^n(Z) \arrow[rd, "\delta^n"] & \\
		S \arrow[twoheadrightarrow, ru, "w"] \arrow[rd, "u"'] & & \Hm^{n+1}(X) \\
		& \Ker\left( d_X^{n+1} \right) \arrow[twoheadrightarrow, ru] &
	\end{tikzcd} \quad \text{komutatif}; \end{equation}
	karena $w$ diketahui merupakan epimorfisme, diagram ini juga menentukan
	$\delta^n$ secara unik.

	Untuk membandingkannya dengan kerucut pemetaan, verifikasi langsung dari
	definisi bahwa diagram berikut komutatif:
% segment-id: o014.li2.chapter3.mapping-cone-and-long-exact-sequences.g011
	\[\begin{tikzcd}
		\Ker(d_X^{n+1}) \oplus \Coker(d_Y^{n-1}) \arrow[d, "{(b,a)}"'] \arrow[hookrightarrow, r] & X^{n+1} \oplus \Coker\left( d_Y^{n-1} \right) \arrow[twoheadrightarrow, r] & \Coker\left( d_{\Cone(f)}^{n-1} \right) \arrow[d, "{d_{\Cone(f)}^n}"] \\
		\Ker(d_Y^{n+1}) \arrow[r, "\sim"] & 0 \oplus \Ker\left( d_Y^{n+1} \right) \arrow[hookrightarrow, r] & \Ker\left( d_{\Cone(f)}^{n+1} \right)
	\end{tikzcd}\]
	Definisi panah-panah horizontal seharusnya jelas. Karena bagian berbentuk
	kipas dalam~\eqref{eqn:connecting-eta-delta-0} komutatif, diagram di atas
	menyiratkan bahwa komposisi
% segment-id: o014.li2.chapter3.mapping-cone-and-long-exact-sequences.q009
	\[
	\begin{multlined}
		S \xrightarrow{(-u, v)}
		\Ker\left( d_X^{n+1} \right) \oplus \Coker\left( d_Y^{n-1} \right)
		\\
		{}\xrightarrow[\text{baris pertama}]{\text{diagram di atas}}
		\Coker\left( d_{\Cone(f)}^{n-1} \right)
		\xrightarrow{d_{\Cone(f)}^n}
		\Ker\left( d_{\Cone(f)}^{n+1} \right)
	\end{multlined}
	\]
	ialah $0$. Oleh karena itu,
	$S \xrightarrow{(-u, v)} \Ker\left( d_X^{n+1} \right) \oplus
	\Coker\left( d_Y^{n-1} \right) \to
	\Coker\left( d_{\Cone(f)}^{n-1} \right)$ berfaktor secara unik sebagai
	$S \to \Hm^n(\Cone(f)) \hookrightarrow
	\Coker\left( d_{\Cone(f)}^{n-1} \right)$.
	Kita nyatakan bahwa diagram berikut komutatif:
% segment-id: o014.li2.chapter3.mapping-cone-and-long-exact-sequences.g012
	\begin{equation}\label{eqn:connecting-eta-delta-1}\begin{tikzcd}
		S \arrow[ddd, bend right=60, "w"'] \arrow[dd] \arrow[r, "{(-u, v)}"] & \Ker\left( d_X^{n+1} \right) \oplus \Coker\left( d_Y^{n-1} \right) \arrow[dd, "\text{jelas}"] \arrow[r, "\text{proyeksi}"] & \Ker\left( d_X^{n+1} \right) \arrow[twoheadrightarrow, d] \\
		& & \Hm^{n+1}(X) \arrow[hookrightarrow, d] \\
		\Hm^n\left( \Cone(f) \right) \arrow[d, "\simeq"', "{\Hm^n(\Phi)}"] \arrow[hookrightarrow, r] & \Coker\left( d_{\Cone(f)}^{n-1} \right) \arrow[r, "\text{diinduksi oleh $\beta(f)$}"' inner sep=0.5em] \arrow[d, "\text{diinduksi oleh $\Phi$}"' inner sep=0.2em] & \Coker\left( d_X^n \right) \\
		\Hm^n(Z) \arrow[hookrightarrow, r] & \Coker\left( d_Z^{n-1} \right) . &
	\end{tikzcd}\end{equation}
	Memang, komutativitas setiap persegi dapat diverifikasi secara rutin.
	Berdasarkan persegi kanan atas dalam~\eqref{eqn:connecting-eta-delta-0}
	dan $\Phi=(0,g)$, bagian
% segment-id: o014.li2.chapter3.mapping-cone-and-long-exact-sequences.g013
	\[\begin{tikzcd}
		S \arrow[r, "{(-u, v)}"]  \arrow[bend right=60, d, "w"'] & \Ker\left( d_X^{n+1} \right) \oplus \Coker\left( d_Y^{n-1} \right) \arrow[d] \\
		\Hm^n(Z) \arrow[hookrightarrow, r] & \Coker\left( d_Z^{n-1} \right)
	\end{tikzcd}\]
	dari diagram~\eqref{eqn:connecting-eta-delta-1} juga komutatif. Selain itu,
	karena $\Hm^n(Z)\to\Coker(d_Z^{n-1})$ merupakan monomorfisme, komposisi
	dengan panah ini menunjukkan bahwa bagian lengkung kiri yang tersisa
	dalam~\eqref{eqn:connecting-eta-delta-1} juga komutatif.

	Perhatikan bahwa kedua komposisi
	\[
	\begin{aligned}
		\Hm^n(\Cone(f)) &\hookrightarrow
		\Coker \left( d_{\Cone(f)}^{n-1} \right)
		\to \Coker(d_X^n), \\
		\Hm^n(\Cone(f)) &\xrightarrow{\Hm^n(\beta(f))}
		\Hm^{n+1}(X) \hookrightarrow \Coker(d_X^n)
	\end{aligned}
	\]
	sama, karena keduanya diinduksi oleh proyeksi
	ke $X^{n+1}$. Dengan demikian,
	\eqref{eqn:connecting-eta-delta-1} selanjutnya memberikan diagram komutatif
% segment-id: o014.li2.chapter3.mapping-cone-and-long-exact-sequences.g014
	\[\begin{tikzcd}
		S \arrow[d, "w"'] \arrow[rr, "-u"] & & \Ker\left( d_X^{n+1} \right) \arrow[d] \\
		\Hm^n(Z) \arrow[r, "\sim", "{\Hm^n(\Phi)^{-1}}"' inner sep=0.5em] & \Hm^n\left( \Cone(f) \right) \arrow[r, "{\Hm^n(\beta(f))}"' inner sep=0.5em] & \Hm^{n+1}(X) .
	\end{tikzcd}\]
	Berdasarkan~\eqref{eqn:connecting-eta-delta-0.5}, ini menyiratkan bahwa
	komposisi $S\xrightarrow{w}\Hm^n(Z)\xrightarrow{\delta^n}\Hm^{n+1}(X)$
	dan $S\xrightarrow{w}\Hm^n(Z)\xrightarrow{\eta^n}\Hm^{n+1}(X)$
	saling berlawanan tanda. Karena $w$ merupakan epimorfisme, diperoleh
	$\eta^n=-\delta^n$. Terbukti.
\end{bukti}

% segment-id: o014.li2.chapter3.mapping-cone-and-long-exact-sequences.co001
\begin{korolari}\label{prop:connecting-eta-delta-2}
	Dalam situasi di atas, $(\eta')^n=\delta^n$ untuk setiap $n\in\Z$.
\end{korolari}

% segment-id: o014.li2.chapter3.mapping-cone-and-long-exact-sequences.pf004
\begin{bukti}
	Gunakan notasi Lema~\ref{prop:cone-qis}. Berdasarkan
	Proposisi~\ref{prop:connecting-eta-delta}, cukup dibuktikan bahwa
	$(\eta')^n+\eta^n=0$. Namun, persamaan terakhir merupakan akibat
	langsung dari fakta bahwa
	$\alpha(g)\Phi+\Phi'\beta(f):\Cone(f)\to\Cone(g)$ homotopik nol
	(Lema~\ref{prop:cone-qis} (ii)).
\end{bukti}

% segment-id: o014.li2.chapter3.mapping-cone-and-long-exact-sequences.co002
\begin{korolari}\label{prop:cone-connecting}
	Untuk sembarang morfisme $f:X\to Y$ dalam $\cate{C}(\mathcal{A})$,
	morfisme penghubung
	$\xi^n:\Hm^{n+1}(X)\simeq\Hm^n(X[1])\to\Hm^{n+1}(Y)$ dalam barisan
	eksak panjang~\eqref{eqn:cone-long-exact-sequence} sama dengan
	$\Hm^{n+1}(f)$.
\end{korolari}

% segment-id: o014.li2.chapter3.mapping-cone-and-long-exact-sequences.pf005
\begin{bukti}
	Terapkan Lema~\ref{prop:cone-qis} pada barisan eksak pendek
	$0\to Y\to\Cone(f)\to X[1]\to0$. Diperoleh kuasi-isomorfisme
	$(0,\beta(f))$ dan diagram komutatif dengan baris-baris eksak
% segment-id: o014.li2.chapter3.mapping-cone-and-long-exact-sequences.g015
	\[\begin{tikzcd}
		0 \arrow[r] & Y \arrow[r, "\alpha(f)"] & \Cone(f) \arrow[r, "{\alpha(\alpha(f))}"] & \Cone(\alpha(f)) \arrow[d, "{\Phi = (0, \beta(f))}"] \arrow[r] & 0 \\
		0 \arrow[r] & Y \arrow[r, "\alpha(f)"'] \arrow[equal, u] & \Cone(f) \arrow[r, "\beta(f)"'] \arrow[equal, u] & X[1] \arrow[r] & 0
	\end{tikzcd}\]
	Morfisme $\psi:\Cone(\alpha(f))\to X[1]$ yang dibahas dalam
	Lema~\ref{prop:cone-alpha} tepat merupakan $(0,\beta(f))$ di sini:
	komponen derajat $n$ dari morfisme ini ialah proyeksi dari
	$Y^{n+1}\oplus\Cone(f)^n=Y^{n+1}\oplus X^{n+1}\oplus Y^n$ ke
	$X^{n+1}$. Diagram komutatif pada lema tersebut dengan demikian memberikan
% segment-id: o014.li2.chapter3.mapping-cone-and-long-exact-sequences.q010
	\begin{equation}\label{eqn:cone-connecting-aux}
		\text{dalam $\cate{K}(\mathcal{A})$} \quad f[1] \circ (0, \beta(f)) + \beta(\alpha(f)) = 0 .
	\end{equation}

	Sekarang, ingat bahwa
	\[
		\xi^n:\Hm^n(X[1])\to\Hm^{n+1}(Y)
	\]
	merupakan morfisme penghubung dari barisan eksak pendek berikut:
% segment-id: o014.li2.chapter3.mapping-cone-and-long-exact-sequences.q011
	\[ 0 \to Y \xrightarrow{\alpha(f)} \Cone(f) \xrightarrow{\beta(f)} X[1] \to 0, \]
	dan penerapan Proposisi~\ref{prop:connecting-eta-delta} pada barisan eksak
	pendek ini segera menunjukkan bahwa morfisme komposit
% segment-id: o014.li2.chapter3.mapping-cone-and-long-exact-sequences.q012
	\[ \Hm^n(\Cone(\alpha(f))) \xrightarrow[\sim]{\Hm^n((0, \beta(f)))} \Hm^n(X[1]) \xrightarrow{\xi^n} \Hm^{n+1}(Y) \]
	sama dengan $-\Hm^n(\beta(\alpha(f)))$. Akan tetapi,
	\eqref{eqn:cone-connecting-aux} menyiratkan
	$\Hm^n(\beta(\alpha(f)))=-\Hm^n(f[1])\circ
	\Hm^n((0,\beta(f)))$. Jadi,
	$\xi^n=\Hm^n(f[1])=\Hm^{n+1}(f)$.
\end{bukti}

% segment-id: o014.li2.chapter3.mapping-cone-and-long-exact-sequences.co003
\begin{korolari}\label{prop:quism-cone}
	Dengan notasi di atas, $f$ merupakan kuasi-isomorfisme jika dan hanya jika
	$\Hm^n(\Cone(f))=0$ untuk setiap $n\in\Z$.
\end{korolari}

% segment-id: o014.li2.chapter3.mapping-cone-and-long-exact-sequences.pf006
\begin{bukti}
	Substitusikan Korolari~\ref{prop:cone-connecting} ke dalam barisan eksak
	panjang~\eqref{eqn:cone-long-exact-sequence}; kesimpulannya langsung
	terlihat.
\end{bukti}
